\centerline{\bf \Huge Графы I}
\centerline{\bf \Huge Уровень: шкибиди туалет.}

\problem{1}
Каждый член партии доверяет пяти однопартийцам, но никакие двое не доверяют друг другу. При каком минимальном размере партии такое возможно?

\problem{2}
В стране из каждого города выходит не более 4   дорог. Также известно, что между любыми двумя городами существует путь, проходящий не более чем по двум другим городам. Докажите, что в этой стране не более 53   городов.

\problem{3}
В клетчатом квадрате $10 \times 10$, вначале пустом, Жанна закрашивает по одной клетке, вписывая в каждую только что закрашенную клетку количество граничащих с нею (по стороне) ранее закрашенных клеток. Чему будет равна сумма всех чисел в квадрате, когда будут закрашены все клетки?

\problem{4}
В графе степень каждой вершины не меньше $k \geqslant 2$. Докажите, что в этом графе существует простой цикл длины не меньше $k+1$.

\problem{5}
В компании из ста тысяч человек среди любых десяти есть трое попарно знакомых. Докажите, что можно выбрать восьмерых из них так, чтобы любой из оставшихся был знаком с кем-то из этих восьмерых.

\problem{6}
В стране 20 городов, причем между любыми двумя городами проложена дорога. Министерство путей сообщения может закрыть на ремонт любую дорогу из четырех, образующих циклический маршрут. Может ли после нескольких таких операций остаться только 19 дорог?